\section{提问备查（Q\&A）}

以下是教授可能提问的预设回答。每条的措辞都面向「快速、准确、不夸大」。

\begin{qanda}{你为什么相信自己算对了？}
三条相互独立的物理判据，任何一条不过都不放过：\textbf{① 多近似方法对比}（Table 2 体积分 vs Mie 级数求和，两条独立数学路径四通道 $<0.03\%$；Grahn 双路径 + miepython + 独立远场投影四个来源互证；双盘用有限元 900G 与边界元 128G 两种求解器）。\textbf{② 极限退化}（Table 2 在 $kr\to0$ 逐条退化回 Table 1，$<1\%$；Mie 退化回 Rayleigh $x^4$；$Q_{\rm ext}\to2$）。\textbf{③ 能量守恒与光学定理}（$C_{\rm ext}=C_{\rm sca}+C_{\rm abs}$ 到 $10^{-10}$；光学定理 $S(0)$ 双路差 $10^{-14}$；Grahn 消光/散射两式差 $10^{-7}$；双盘功率闭合 balance $=1.0$）。三条叠加排除了「实现 bug」「公式抄错」「物理不自洽」三类错误。
\end{qanda}

\begin{qanda}{MD/EQ 为什么有 $\sim2\%$ 网格残差？}
双金盘在 $x=0.36$ 是磁共振峰，表面附近场梯度极大，多极矩体积分对表面四面体分辨率高度敏感。这是表面场奇异性主导的收敛行为，Richardson 外推显示收敛阶 $p\approx1.3$（一阶收敛，h 减半误差只减约 40\%）。$0.3$ 网格（1070 万单元）是当前集群 900G 内存内可达到的最细网格；继续细化到 $0.2$ 需约 2TB 内存，超集群上限。所以用 Richardson 外推定量估计了收敛值（MD $1.54$--$1.58\times10^{-12}$），并如实披露 $2\%$ 残差，而不是假装它收敛了。
\end{qanda}

\begin{qanda}{为什么用 FEM 900G 不用别的数值方法？}
有三个工程理由：\textbf{① 现成工具链}——COMSOL 6.3 的 Java API 与集群已有镜像/Parasolid 库已打通，不用另装别的求解器；\textbf{② 论文对齐}——Alaee 2018 自己用 COMSOL 做球体验证，沿用同一软件便于方法对齐；\textbf{③ 集群够大}——magnus 更新后可申请到 900G。但也验证了更省的替代：\textbf{边界元法用 128G 就跑通了同一几何}（省约 7 倍），这是后续深化网格收敛的正确方向（表面型方法只离散表面，内存不随体积膨胀）。
\end{qanda}

\begin{qanda}{光学定理怎么验证的？}
两条独立路径：\textbf{① 级数求和}——把 Mie 系数代入消光截面级数，直接求和；\textbf{② 前向极限}——由光学定理 $C_{\rm ext}=\frac{4\pi}{k^2}\Re[S(0)]$，从前向散射振幅的虚部取极限。两条路径复 $S(0)$ 差 $1.2\times10^{-14}$。此外 Grahn 的 Eq.~(22)（消光）与 Eq.~(20)（散射）在 $x=0.5$ 处相对差 $5.96\times10^{-7}$，也间接验证了消光/散射的物理关系。
\end{qanda}

\begin{qanda}{toroidal 多极为什么不是独立自由度？}
这是 Fernandez-Corbaton 2017 的严格证明（\emph{Sci. Rep.} 7, 7527）：把电偶极的精确展开按电部分与 toroidal 部分拆分后，\textbf{两部分各自含有不可观测的非辐射（off-shell）分量}，只有两者求和才抵消。因此 toroidal 部分无法通过辐射或外场耦合被单独测量，它只是电偶极系数对电磁尺寸（$kr$）展开的次低阶修正项（对应 Table 1 的 $k^2/10$ 项）。推导并退化验证了这一结论。
\end{qanda}

\begin{qanda}{Table 1 和 Table 2 是什么关系？}
Table 1 是\textbf{长波近似}（$kr\to0$ 截断），只保留多极矩的最低阶；Table 2 是\textbf{精确表达式}，保留球 Bessel 核 $j_l(kr)$ 的全部阶。两者关系是 Table 1 = Table 2 的 $kr\to0$ 极限（$j_l/(kr)^l\to1/(2l+1)!!$）。逐条验证了 ED/MD/EQ/MQ 四条退化（$<1\%$）。Table 2 在任意尺寸参数下都精确（Alaee 用 Mie 验证与精确解完全一致）。
\end{qanda}

\begin{qanda}{你复现的误差有多大？}
球体（Fig.~1/2）的 Table 2 vs Mie 四通道最大误差：介电球 $<0.21\%$，金球（加密网格）$<0.03\%$。Grahn 双路径映射：长波路径 $5.7\times10^{-7}$，有限核路径 $2.9\times10^{-4}$。双盘（Fig.~3）没有解析真值，用两种数值方法（FEM/BEM）互验，并给出 $2\%$ 的网格残差上界。这些误差都是「方法与解析真值的偏差」，不是「与论文图的视觉差」。
\end{qanda}

\begin{qanda}{这个复现的意义是什么？}
一是\textbf{验证了多极展开理论链}（Grahn 映射 $\to$ FC 精确偶极 $\to$ Alaee 精确多极矩）在数值上自洽，球体上与方法无关的 Lorenz--Mie 解一致到 $10^{-4}$ 量级；二是\textbf{打通了非球形（双盘）的频域真解}，用两种数值求解器交叉验证；三是\textbf{沉淀了可审计的方法链}——从电流分布到多极矩、散射截面，每一步都有独立物理验证，不是「画图看起来像」。
\end{qanda}

\begin{qanda}{你用的材料数据是什么？}
金球用 Johnson--Christy 1972 的复折射率（$n,k$ 随波长，支持到 1935 nm）；介电球用 $\varepsilon_r=2.5^2=6.25$。双盘的 COMSOL 也用 Johnson--Christy 金数据（色散插值，禁止外推）。材料数据源都做了 provenance 记录。
\end{qanda}

\begin{qanda}{最容易错的地方在哪？你怎么避免？}
多极矩公式的\textbf{系数}（如 ED 的 $3$、MQ 的 $15$、toroidal 的 $k^2/10$）最易抄错。避免的方法是\textbf{退化检查}：每个系数都要求它在 $kr\to0$ 退化回长波近似的已知结果，系数由退化固定而不是「照抄论文」。例如 MQ 的 $15\times j_2/(kr)^2\to15\times1/15=1$，任何一步系数错都会在退化里暴露。
\end{qanda}
